\subsection{An ordinary obstruction to the large-prime quotient test}
\label{sec:b3-ordinary-obstruction}

\paragraph{Finding A23 (the large-prime quotient test can fail over every extension).}
Let $C/\mathbf F_4$ be the smooth projective model of
\[
                   y^2+y=x+x^{-1}+(x-1)^{-1}.
 \tag{B3O.1}
\]
It is an ordinary genus-two curve.  Put
$Q=4^m$, $N_m=\#C(\mathbf F_Q)$ and
$J_m=\#J_C(\mathbf F_Q)$ for $m\ge1$.
For every prime $\ell\mid J_m$ one has
\[
                              \ell N_m<2J_m.
 \tag{B3O.2}
\]
Consequently the sufficient conditions in recovered A1
B.3, formula (B.CL.2), can never hold for a prime quotient
of this actual Jacobian over any constant extension of
$\mathbf F_4$, even if the primitive-root requirement is omitted.
Nevertheless the Raynaud theta support of this same curve has
nonzero points of exact order $4159$ over $\mathbf F_{4^7}$,
and $\langle2,-1\rangle=(\mathbf Z/4159)^\times$.
Thus (B3O.2) is a failure of a sufficient numerical test, not
an absence of geometric theta witnesses.

\begin{proof}
The equation is defined over $\mathbf F_2$.  Its three simple
poles at $0,1,\infty$ give genus two and $2$-rank two by the
Artin--Schreier formulas \cite[Lemma~2.6]{PriesZhu2012}.
There are three rational points over each of $\mathbf F_2$
and $\mathbf F_4$: the three ramified points account for all
$\mathbf F_2$-points; if $u^2+u+1=0$, the remaining
$\mathbf F_4$-values of $x$ are $u,u^2$, and the right side
of (B3O.1) has respective values $u^2,u$, both of absolute
trace one.  These have no Artin--Schreier preimages.
These are also the exact inputs (LW.12) of A13.

Let $\alpha_1,\ldots,\alpha_4$ be the Frobenius eigenvalues
over $\mathbf F_2$.  The two point counts imply
$\sum_i\alpha_i=0$ and $\sum_i\alpha_i^2=2$.
The genus-two functional equation therefore yields the
characteristic polynomial
\[
                         P_2(T)=T^4-T^2+4,
 \tag{B3O.3}
\]
as in (LW.13).  If $\beta_1,\beta_2$ are the roots of
$T^2-T+4$, the four squares $\alpha_i^2$ are
$\beta_1,\beta_1,\beta_2,\beta_2$.
Thus the Frobenius polynomial over $\mathbf F_4$ is
$(T^2-T+4)^2$; no isogeny classification is needed here.
Define
\[
 t_m=\beta_1^m+\beta_2^m,\qquad e_m=Q+1-t_m.
 \tag{B3O.4}
\]
The recurrence $t_0=2$, $t_1=1$ and
$t_m=t_{m-1}-4t_{m-2}$ shows that $t_m,e_m$ are integers.
The curve and Jacobian Frobenius point-count formulas give
\[
 \begin{aligned}
 N_m&=Q+1-2t_m=2e_m-(Q+1),\\
 J_m&=(1-\beta_1^m)^2(1-\beta_2^m)^2=e_m^2.
 \end{aligned}
 \tag{B3O.5}
\]
For the Jacobian formula see
\cite[Theorem~19.1(b)]{MilneAV}.
Since $|\beta_i|=2$, the Weil estimate gives
$e_m\ge(\sqrt Q-1)^2>0$.

If a prime $\ell$ divides $J_m=e_m^2$, then
$\ell\mid e_m$ and $\ell\le e_m$.
Also $N_m\ge0$, since it is a point count.  Hence
\[
 \ell N_m\le e_mN_m
       =2e_m^2-(Q+1)e_m
       <2e_m^2=2J_m,
 \tag{B3O.6}
\]
which proves (B3O.2) for every $m$ and every such prime.

The recovered source \cite[B.3, formula~(B.CL.2)]{A1run}
requires, for a surjection
$J_C(\mathbf F_Q)\twoheadrightarrow\mathbf Z/\ell$
with kernel $K$, both $N_m\ge2|K|$ and $\ell>N_m$,
in addition to $2$ being primitive modulo $\ell$.
Since $|K|=J_m/\ell$, its first inequality is exactly
$\ell N_m\ge2J_m$, contradicted by (B3O.2).
No primitive-divisor theorem or large-prime-factor theorem
can universally force these sufficient inequalities:
they are false for this ordinary curve over every extension
of its stated base field.

Finally A13 proves that the same model over $\mathbf F_2$
has at least four Raynaud-theta points of exact order
$4159$ over $\mathbf F_{128}=\mathbf F_{2^7}$, and that
$\langle2,-1\rangle=(\mathbf Z/4159)^\times$.
The inclusion
$\mathbf F_{2^7}\subset\mathbf F_{2^{14}}=\mathbf F_{4^7}$
preserves their order and theta incidence.  They therefore
remain valid geometric witnesses on the curve (B3O.1)
over its specified base $\mathbf F_4$.
The primitive-root condition in B.3 is stronger than this
signed-doubling condition: A13 has
$\operatorname{ord}_{4159}(2)=2079$, not $4158$.
The two arguments are consistent because B.3 gives only
one sufficient test and asserts no converse.
\end{proof}

\paragraph{Exact source locator for A23.}
The recovered B.3 is the complete claim and proof titled
\emph{The corrected pair estimate and its precise limitation}
in \path{ANABELIAN_A1_paper.tex}, lines 2488--2546 of the
retained checkpoint-03 text.  A13 supplies formulas
(LW.12)--(LW.17) and its complete theta-incidence proof.
The all-extension inequality (B3O.6) is symbolic; the
accompanying finite recurrence checks are arithmetic
cross-checks, not the proof of its universal quantifier.
